\chapter{fixed60：收益、成本、entry/exit 是怎么算的}

\section{本章要解决的困惑}

entry 以后，fast 回测怎么知道什么时候退出？当前默认最简单路径是 fixed60：

\[
exit\_ts = entry\_ts + 60s.
\]

这不是最优 exit，也不是 path prediction，只是固定持有 60 秒的读数。它把复杂问题压成一个可重复的标签。

\section{mid 口径}

如果用 \code{fixed60_mid}，收益是：

\[
Y^{mid}_{60}
= direction \times 10000 \times \log\frac{mid_{t+60s}}{mid_t}.
\]

其中：

\[
direction =
\begin{cases}
1, & long/buy,\\
-1, & short/sell.
\end{cases}
\]

例如 long：

\[
mid_t=0.1122,\quad mid_{t+60}=0.1128.
\]

\[
Y^{mid}_{60}
=10000 \times \log(0.1128/0.1122)
\approx 53.3\text{ bps}.
\]

如果 short，同样的价格上涨会是负收益。

\section{taker 可执行口径}

如果用 \code{fixed60_taker} 或 top-of-book taker 近似，不能只看 mid。因为做多入场打 ask，退出卖出打 bid：

\[
Y^{exec,long}_{60}
= 10000 \times \log\frac{bid_{t+60s}}{ask_t}.
\]

做空入场卖 bid，退出买 ask：

\[
Y^{exec,short}_{60}
= 10000 \times \log\frac{bid_t}{ask_{t+60s}}.
\]

这就是为什么零手续费也会有执行成本。mid 正收益可能被 entry/exit spread 吃掉。

\section{小订单簿例子}

入场时：

\begin{lstlisting}
bid_t = 0.1120
ask_t = 0.1124
mid_t = 0.1122
\end{lstlisting}

60 秒后：

\begin{lstlisting}
bid_exit = 0.1126
ask_exit = 0.1130
mid_exit = 0.1128
\end{lstlisting}

long mid 口径：

\[
Y^{mid,long}_{60}
=10000\log(0.1128/0.1122)
\approx 53.3.
\]

long executable：

\[
Y^{exec,long}_{60}
=10000\log(0.1126/0.1124)
\approx 17.8.
\]

同一个路径，mid 看起来赚 53.3 bps，可执行只剩 17.8 bps。差别就是 spread/cross 成本。

\section{fast entries 和 exits 的对应关系}

\code{fast_strategy_backtest.py} 里：

\begin{itemize}[leftmargin=2em]
  \item entry row 记录 \code{entry_mid}、\code{entry_bid}、\code{entry_ask}、\code{actual_exposure}；
  \item exit row 记录 \code{exit_mid}、\code{exit_bid}、\code{exit_ask}、\code{raw_bps_mid_fixed60}、\code{net_bps}；
  \item daily row 汇总每一天 entries、exits、actual exposure、weighted bps；
  \item summary 记录总结果和 profile manifest。
\end{itemize}

注意：\code{entries.parquet} 是所有候选和 capacity 决策的记录；真正有仓位的是：

\begin{lstlisting}
actual_exposure > 0
\end{lstlisting}

\section{fixed60 的局限}

用户指出的最大问题是对的：fixed60 太粗。它没有回答：

\begin{itemize}[leftmargin=2em]
  \item 最大收益是多少；
  \item 最大值什么时候出现；
  \item 峰值后什么时候回到 0bps；
  \item 是否应该 5 秒、20 秒、120 秒退出；
  \item entry 结构是否对应不同 path distribution。
\end{itemize}

所以 fixed60 只是第一版可重复读数，不是完整预测对象。

\section{手动检查点}

给一笔 long entry：

\begin{lstlisting}
entry_bid = 0.1120
entry_ask = 0.1124
exit_bid  = 0.1126
exit_ask  = 0.1130
\end{lstlisting}

先算 mid，再算 executable。最后问：

\begin{quote}
mid 的收益是不是被 top-of-book cross 吃掉了？
\end{quote}

这比只看 \code{net_median} 更接近执行现实。

